\begin{figure}[!htbp]
\centering
\resizebox{0.9\linewidth}{!}{%
\begin{tikzpicture}[font=\small,>=Stealth,line width=0.7pt]
% LUT
\node[draw,fill=black!8,minimum width=2.6cm,minimum height=3cm,align=center] (lut) at (0,0) {Look-up\\table\\(LUT)};
% inputs x1, x3, x4 from the layer below / neighbours
\foreach \yy/\lab in {1.1/x_1,-0.37/x_3,-1.1/x_4}{
  \node[anchor=east] at (-2.7,\yy) {$\lab$};
  \draw (-2.6,\yy) -- (-1.3,\yy);
}
% x2 is the feedback wire
\node at (-2.2,0.62) {$x_2$};
\draw (7.4,0) -- (7.4,-2.6) -- (-3.4,-2.6) -- (-3.4,0.37) -- (-1.3,0.37);
\node[font=\footnotesize,anchor=north] at (2.0,-2.6) {feedback of $y$ through the routing fabric};
% LUT output
\draw (1.3,0) -- (2.3,0);
\fill (2.3,0) circle (2pt);
% bypass to MUX
\draw (2.3,0) -- (2.3,0.6) -- (6.0,0.6);
% FF
\draw (2.3,0) -- (2.3,-0.6) -- (3.2,-0.6);
\node[draw,minimum width=1.6cm,minimum height=1.8cm] (ff) at (4.0,-0.6) {FF};
\node[anchor=west,font=\footnotesize] at (3.2,-0.3) {D};
\node[anchor=east,font=\footnotesize] at (4.8,-0.3) {Q};
\draw (4.8,-0.6) -- (6.0,-0.6);
% clock
\draw (2.9,-2.1) -- (2.9,-1.2) -- (3.2,-1.2);
\draw (3.2,-1.05) -- (3.4,-1.2) -- (3.2,-1.35);
\node[font=\footnotesize,anchor=north] at (2.9,-2.1) {clk};
% MUX
\draw (6.0,1.0) -- (6.0,-1.0) -- (6.8,-0.6) -- (6.8,0.6) -- cycle;
\node[rotate=90,font=\footnotesize] at (6.4,0) {MUX};
\draw (6.8,0) -- (7.4,0);
\fill (7.4,0) circle (2pt);
\draw[->] (7.4,0) -- (8.3,0) node[right] {$y$};
\end{tikzpicture}}
\caption{Logic block used as a recurrent unit (schematic redrawn after the generic LUT--FF--MUX logic cell). The MUX selects the registered path $Q$; the block output $y$ is routed back to LUT input $x_2$, so the LUT sees the new input bit $x_1$ together with the unit's own previous output.}
\label{fig:logic-block}
\end{figure}
